\subsection{The matrix Fisher operator}
\label{sec:fisher-operator}

The matrix Fisher operator separates diagonal scale, symmetric off-diagonal,
and skew-symmetric rotational directions.  Its rotational eigenvalue is
\(J-1\), which vanishes exactly in the Gaussian case.

Let $f$ be a probability density and write $g=\sqrt f$.
In this subsection we assume
\begin{equation}
 \begin{gathered}
 f>0\quad\text{Lebesgue-almost everywhere},\qquad
 g\in W^{1,2}(\R),\qquad xg'(x)\in L^2(\R),\\
 \E X=0,\qquad \E X^2=1,
 \qquad X\sim f(x)\,dx.
 \end{gathered}
 \label{eq:fisher-regularity}
\end{equation}
The derivatives in these assumptions are weak derivatives.
Define the location score $s=-f'/f=-2g'/g$ almost everywhere.  Then
\begin{equation}
 \E s(X)=0,\qquad \E[Xs(X)]=1,\qquad
 J:=\E s(X)^2=4\|g'\|_2^2.
 \label{eq:score-moments}
\end{equation}
Define the scale score and its second moment by
\begin{equation}
 d(x):=xs(x)-1,\qquad
 \kappa:=\E d(X)^2=4\|xg'\|_2^2-1.
 \label{eq:kappa}
\end{equation}
The two integration-by-parts identities in
\eqref{eq:score-moments} follow from these assumptions.
Indeed,
$f'=2gg'$ belongs to $L^1$, and $xf'$ belongs to $L^1$ by
Cauchy--Schwarz and the finite second moment.  Choose a smooth compactly
supported cutoff $\chi$ equal to one near zero.  Applying integration by
parts first to $\chi(x/L)$ and then to
$x\chi(x/L)$ gives the identities as $L\to\infty$.

For independent copies $X_1,\ldots,X_n$ of $X$, write
$\mathbf X^{(n)}=(X_1,\ldots,X_n)\sim\mu_n$.  For
$A\in M_n(\R)$, define the score of the pushforward by $I+tA$ at $t=0$ by
\begin{equation}
 S_A(\mathbf X^{(n)})=\sum_{i,j=1}^n A_{ij}
       \bigl(s(X_i)X_j-\delta_{ij}\bigr)
       =\sum_i A_{ii}d(X_i)+\sum_{i\ne j}A_{ij}s(X_i)X_j.
 \label{eq:matrix-score}
\end{equation}

\begin{theorem}[The matrix Fisher operator]
\label{thm:fisher-spectrum}
Under \eqref{eq:fisher-regularity}, let \(s\) and \(d\) be the location and
scale scores defined above, and define the matrix Fisher operator
$\mathcal I$ by
$\langle A,\mathcal I B\rangle_F=\E[S_AS_B]$.
Then
\begin{equation}
 \boxed{\mathcal I A=JA+A^*
                     +(\kappa-J-1)\diag A.}
 \label{eq:fisher-operator-spectrum}
\end{equation}
Here $\diag A$ is the matrix with diagonal entries $A_{ii}$ and zero
off-diagonal entries.  The diagonal, symmetric off-diagonal, and skew-symmetric matrix subspaces
are orthogonal for the Frobenius inner product and have eigenvalues
$\kappa$, $J+1$, and $J-1$, respectively.  Equivalently,
\begin{equation}
 \begin{aligned}
 \langle A,\mathcal I A\rangle_F
 &=\kappa\|\diag A\|_F^2
   +(J+1)\left\|\frac{A+A^*}2-\diag A\right\|_F^2\\
 &\quad +(J-1)\left\|\frac{A-A^*}2\right\|_F^2.
 \end{aligned}
 \label{eq:fisher-decomposition}
\end{equation}
The eigenvalue multiplicities are, respectively,
\begin{equation}
 n,\qquad n(n-1)/2,\qquad n(n-1)/2.
 \label{eq:fisher-multiplicities}
\end{equation}
Moreover, $\kappa>0$ and $J\ge1$.
Here $J=1$ holds precisely for the
standard Gaussian density.  Thus, for a non-Gaussian $f$,
$\min\{\kappa,J-1\}>0$ and
\begin{equation}
 \begin{gathered}
 \min\{\kappa,J-1\}\|A\|_F^2\le\langle A,\mathcal I A\rangle_F
 \le C_f\|A\|_F^2,\\
 C_f:=\max\{\kappa,J+1\}.
 \end{gathered}
 \label{eq:fisher-ellipticity}
\end{equation}
\end{theorem}

\begin{proof}
The diagonal variables $d(X_i)$ are independent and centered.
Their covariances with every off-diagonal term $s(X_j)X_k$, $j\ne k$,
vanish.
After any shared coordinate is integrated out, the remaining factor has
expectation \(\E X\) or \(\E s(X)\).
Two off-diagonal terms can have a nonzero covariance only when their
ordered index pairs coincide or are reversed.  Those covariances are,
respectively, $J$ and $(\E Xs(X))^2=1$.  Consequently
\begin{equation}
 \E S_A^2
 =\kappa\sum_i A_{ii}^2+
   \sum_{i<j}\{J(A_{ij}^2+A_{ji}^2)+2A_{ij}A_{ji}\}.
 \label{eq:fisher-pair-form}
\end{equation}
Resolving each pair into its symmetric and skew-symmetric parts proves
the decomposition.  In particular, the calculation uses centering and
independence and imposes no parity condition on $f$.

Cauchy--Schwarz applied to $\E Xs(X)=1$ gives $J\ge1$.
Equality implies $s(X)=X$ almost surely, hence
$g'=-xg/2$ almost everywhere and $f$ is the standard Gaussian.
If $\kappa=0$, then $2xg'+g=0$ almost everywhere on each of the two
half-lines.  The absolutely continuous solutions there are constant
multiples of $|x|^{-1/2}$, which cannot give a positive probability
density.  This proves $\kappa>0$ and the asserted bounds.
\end{proof}

\subsection{Hellinger paths, polar decomposition, and symmetric perturbations}
\label{sec:hellinger-paths}

For orientation, write the left polar decomposition of an invertible matrix
as \(T=BO\), where \(B=(TT^*)^{1/2}\) is positive and \(O\) is
orthogonal.  The covariance defect is \(TT^*-I=B^2-I\), so a
finite second moment detects the symmetric factor.  The skew tangent of the
orthogonal factor has Fisher eigenvalue \(J-1\), which is positive precisely
off the Gaussian law.  The full matrix Fisher operator decomposition above packages
these two mechanisms without requiring a nonlinear separation of \(B\) and
\(O\) at every step.

For $A\in\GL(n,\R)$ define a unitary operator on
$L^2(\R^n,dx)$ by
\[
 (U_Av)(y)=|\det A|^{-1/2}v(A^{-1}y),
 \qquad
 \Omega_n(x)=\prod_{i=1}^n g(x_i).
\]
Thus $U_A\Omega_n$ is the square root of the Lebesgue density of
$A_\#\mu_n$.  For invertible $A,B$,
\[
 h(A_\#\mu_n,B_\#\mu_n)=\|U_A\Omega_n-U_B\Omega_n\|_2.
\]
For a fixed real matrix $A$, the generator of $U_{\exp(tA)}$ is
\begin{equation}
 G_Av=-\tfrac12\operatorname{tr}(A)v-(Ax)\cdot\nabla v,
 \qquad G_A\Omega_n=\tfrac12S_A\Omega_n.
 \label{eq:half-density-generator}
\end{equation}
These identities hold in $L^2$.
To justify them under \eqref{eq:fisher-regularity}, observe that
$x_j\partial_i\Omega_n\in L^2$ for all $i,j$: the diagonal case
uses $xg'\in L^2$, and the off-diagonal case uses $g'\in L^2$
and $\E X^2=1$.  Cutoff and mollification approximate
$\Omega_n$ in the corresponding weighted Sobolev graph norm, where
the usual change-of-variables differentiation is valid.

\begin{lemma}[A dimension-free Hellinger path bound]
\label{lem:finite-hellinger-upper}
If $A_t$, $0\le t\le1$, is a continuously differentiable path in
$\GL(n,\R)$, then $t\mapsto U_{A_t}\Omega_n$ is continuously
differentiable in $L^2$ and
\begin{equation}
 \begin{aligned}
 \frac d{dt}U_{A_t}\Omega_n
 &=U_{A_t}G_{A_t^{-1}\dot A_t}\Omega_n,\\
 \left\|\frac d{dt}U_{A_t}\Omega_n\right\|_2
 &=\frac12\sqrt{\left\langle A_t^{-1}\dot A_t,
                  \mathcal I(A_t^{-1}\dot A_t)\right\rangle_F}.
 \end{aligned}
 \label{eq:finite-hellinger-path-derivative}
\end{equation}
Consequently,
\begin{equation}
 h((A_1)_\#\mu_n,(A_0)_\#\mu_n)
 \le\frac{\sqrt{C_f}}2
       \int_0^1\|A_t^{-1}\dot A_t\|_F\,dt.
 \label{eq:finite-hellinger-path}
\end{equation}
In particular, if $\|B\|_{\mathrm{op}}\le r<1$, then
\begin{equation}
 h((I+B)_\#\mu_n,\mu_n)
 \le\frac{\sqrt{C_f}}{2(1-r)}\|B\|_F.
 \label{eq:finite-hellinger-upper}
\end{equation}
All constants are independent of $n$.
\end{lemma}

\begin{proof}
The representation law $U_AU_B=U_{AB}$ and
\eqref{eq:half-density-generator} give
\eqref{eq:finite-hellinger-path-derivative}.
Integrate the derivative and use the Fisher upper bound.  For the last
assertion take $A_t=I+tB$ and use
$\|A_t^{-1}\|_{\mathrm{op}}\le(1-r)^{-1}$.
\end{proof}

\subsection{Finite real Fourier sums with dual moments}
\label{sec:fourier-dual-moments}

Assume in addition that \(f\) is non-Gaussian.  The score \(S_A\) is the correct
tangent, but it is not a usable dimension-free test after the law moves:
differentiating a score-based test would require regularity and integrability
not assumed here.  We instead need, for every matrix direction, a test built
from bounded one-coordinate functions, with exact duality against the score and
with \(L^2\)-norm bounded by a constant times the Frobenius norm of that
direction, uniformly in the dimension.  The rank-two
obstruction in \cref{lem:fourier-ode} produces the two off-diagonal coordinate
functions; a separate small-frequency argument using the second moment
produces the diagonal coordinate function.  All three are finite real Fourier
sums at frequencies depending on \(f\) but not on the matrix dimension.

\begin{lemma}[Finite real Fourier sums with dual moments]
\label{lem:fourier-dual-moments}
There exist centered finite real Fourier sums $p,q,r$ such that
\begin{equation}
 \begin{array}{c|cc}
       &s&x\\ \hline
 p     &1&0\\
 q     &0&1
 \end{array}
 \qquad\text{and}\qquad \E[r(X)d(X)]=1.
 \label{eq:fourier-dual-moments}
\end{equation}
The entries in the table denote expectations of the products.
For
\begin{equation}
 H_A(\mathbf X^{(n)}):=\sum_i A_{ii}r(X_i)
              +\sum_{i\ne j}A_{ij}p(X_i)q(X_j),
 \label{eq:fourier-functions}
\end{equation}
one has, for all real matrices $A,B$ of the same size,
\begin{equation}
 \E[H_AS_B]=\langle A,B\rangle_F,
 \qquad \|H_A\|_{L^2(\mu_n)}\le C_{\mathrm{dual}}\|A\|_F,
 \label{eq:matched-fourier-pairing}
\end{equation}
where $C_{\mathrm{dual}}$ is independent of $n$.
\end{lemma}

\begin{proof}
On the real vector space of centered finite real Fourier sums
consider the two functionals
\[
 L_1(u)=\E u'(X)=\E[u(X)s(X)],
 \qquad L_2(u)=\E[Xu(X)].
\]
The integration-by-parts identity holds since $u,u'$ are bounded.
Both functionals are unchanged by subtracting a constant.  If this
moment map had rank at most one, all its two-by-two determinants
would vanish, including after complexification.  Writing
$\phi(\xi)=\E e^{i\xi X}$ and testing with the centered
exponentials at frequencies $\xi,\eta$ would then give
\[
 \eta\phi'(\xi)\phi(\eta)
       -\xi\phi(\xi)\phi'(\eta)=0
 \qquad(\xi,\eta\in\R).
\]
Here $L_1(e^{i\xi x})=i\xi\phi(\xi)$ and
$L_2(e^{i\xi x})=-i\phi'(\xi)$.  The global linear ODE
argument in \cref{lem:fourier-ode} would force a centered
Gaussian law, a contradiction.  The map therefore has rank two.
Inverting the moment matrix of two finite real Fourier sums
supplies $p,q$ with the displayed dual
moments.  This construction fixes finitely many frequencies
depending only on $f$.

The diagonal function can be chosen using only the second moment.
For $\xi\ne0$ let
$r_0(x)=\cos(\xi x)-\E\cos(\xi X)$.  Dominated convergence
gives
\[
 \frac{\E[Xr_0'(X)]}{\xi^2}
 =-\E\left[X\frac{\sin(\xi X)}\xi\right]
 \longrightarrow-\E X^2=-1
 \qquad(\xi\to0).
\]
Choose a sufficiently small nonzero $\xi$ and set
$r=r_0/\E[Xr_0'(X)]$.  Cutoff integration by parts gives
$\E[r(X)d(X)]=\E[Xr'(X)]=1$.

All diagonal/off-diagonal cross terms in $\E H_AS_B$ vanish
by centering.  Coincident ordered off-diagonal pairs contribute
$\E[ps]\E[qX]=1$.
Reversed pairs contribute $\E[pX]\E[qs]=0$.
This proves the pairing.
For any centered $a,b\in L^2(f)$, independence gives
\begin{equation}
 \left\|\sum_{i\ne j}A_{ij}a(X_i)b(X_j)\right\|_2
 \le\sqrt2\,\|a\|_2\|b\|_2\|A\|_F.
 \label{eq:quadratic-chaos-bound}
\end{equation}
Together with the independent diagonal terms, this proves the last
assertion.
\end{proof}

The test functions \(H_A\) will be paired with the Hellinger tangent along a matrix
path.  To keep that pairing uniformly controlled after moving away from the
identity, we need one uniform derivative bound for \(H_A\Omega_n\).  A naive
term count grows with \(n\), while a direct H\"older estimate asks for
integrability that is unavailable.  The next lemma isolates this non-formal
step.  Its proof groups terms by their index-coincidence patterns, turning every
surviving coefficient array into a Frobenius-controlled contraction.  In
particular, it uses no fourth moment of \(X\) or higher moment of the score.

\begin{lemma}[Generator bound for the functions $H_A$]
\label{lem:fourier-generator-bound}
Let $p,q,r$ be any fixed centered bounded continuously differentiable
functions with bounded derivatives.  Define $H_A$ by
\eqref{eq:fourier-functions}.  There is a constant $C_{\mathrm{gen}}$,
depending on these functions and on $J,\kappa$, such that
\begin{equation}
 \|G_B(H_A\Omega_n)\|_{L^2(dx)}
 \le C_{\mathrm{gen}}\|A\|_F\|B\|_F
 \qquad(A,B\in M_n(\R),\ n\ge1).
 \label{eq:coordinate-function-generator-bound}
\end{equation}
\end{lemma}

\begin{proof}
The issue is not membership in the generator domain in a fixed dimension, but
uniformity as the dimension grows.  The product rule creates up to four
coordinate factors; grouping them by coincidences rather than counting terms
will keep the constants independent of \(n\).

For fixed $n$, both $H_A$ and its first derivatives are bounded.
Consequently
\[
 x_j\partial_i(H_A\Omega_n)
 =H_Ax_j\partial_i\Omega_n
       +x_j(\partial_iH_A)\Omega_n\in L^2(dx)
\]
by the weighted Sobolev regularity of $\Omega_n$ and
$\E X^2=1$.  The cutoff and mollification argument used for
\eqref{eq:half-density-generator} therefore puts $H_A\Omega_n$
in the domain of $G_B$.
The product rule gives
\begin{equation}
 G_B(H_A\Omega_n)
 =\left\{\tfrac12H_AS_B-(Bx)\cdot\nabla H_A\right\}\Omega_n.
 \label{eq:coordinate-function-generator-product}
\end{equation}
Both terms belong to $L^2$ in each finite dimension.  We prove a
dimension-free bound for each term.

We will repeatedly use the following elementary consequence of
independence.  If $u_1,\ldots,u_m$ are centered $L^2(f)$ functions,
and a coefficient tensor $C$ is indexed by ordered tuples of
\emph{distinct} coordinates, then
\begin{equation}
 \left\|\sum_{i_1,\ldots,i_m}^{\ne}
 C_{i_1\ldots i_m}\prod_{a=1}^m u_a(X_{i_a})\right\|_2
 \le\sqrt{m!}\,\|C\|_{\ell^2}
                     \prod_{a=1}^m\|u_a\|_2.
 \label{eq:canonical-chaos-bound}
\end{equation}
In the square of this norm, a mixed term vanishes unless the two
tuples have the same index set.
There are at most $m!$ matching permutations.
Cauchy--Schwarz for the local moments gives a bound on each index
pattern.
A second Cauchy--Schwarz step for the coefficient arrays proves the
full bound.
We use it only for
$m\le4$.

First consider $H_AS_B$.  For a vector $a$ and a zero-diagonal
matrix $M$, abbreviate
\[
 D_a(u)=\sum_i a_i u(X_i),\qquad
 Q_M(u,v)=\sum_{i\ne j}M_{ij}u(X_i)v(X_j).
\]
The product $H_AS_B$ is the sum of the four products obtained by
multiplying $D_{\mathbf a}(r)+Q_{A_0}(p,q)$
and $D_{\mathbf b}(d)+Q_{B_0}(s,x)$, where
$\mathbf a=(A_{ii})_i$, $\mathbf b=(B_{ii})_i$,
$A_0=A-\diag A$, and
$B_0=B-\diag B$.
In every common coordinate, a local product is a bounded function
times one of $d,s,x$ and hence is in $L^2(f)$.  Write each such
local product as its mean plus its centered part.

For a product $D_a(u)D_b(v)$, distinct indices give a coefficient
tensor bounded by $\|a\|_2\|b\|_2$; equal indices give the
Hadamard vector $(a_ib_i)_i$ and the scalar $\sum_i a_ib_i$.
Their norms have the same bound.  For $D_a(u)Q_M(v,w)$, there
are three index patterns: three distinct indices, equality with
the first index of $Q_M$, and equality with its second index.
The first gives the tensor $a\otimes M$.  In either equality
case, centering the local product gives a matrix obtained by
row or column multiplication of $M$ by $a$, and its mean gives
$M^*a$ or $Ma$.  Each norm is at most
$\|a\|_2\|M\|_F$.  The product $Q_M(u,v)D_a(w)$ has exactly
the same three patterns and bounds.

For clarity, all patterns in a quadratic--quadratic product can
also be enumerated.  Consider
\[
 Q_M(u,v)Q_N(w,z)
 =\sum_{i\ne j,\,k\ne l}M_{ij}N_{kl}
                         u_i v_j w_k z_l,
\]
where $u,v$ are bounded and $w,z$ are in $L^2$.  Four distinct
indices give the restricted tensor $M\otimes N$, of norm at
most $\|M\|_F\|N\|_F$.  There are four patterns with exactly
one shared index: $i=k$, $i=l$, $j=k$, and $j=l$.
For example, the pattern $i=k$ is
\[
 \sum_{i,j,l}^{\ne}M_{ij}N_{il}(uw)_i v_jz_l.
\]
Its centered local product has coefficient norm bounded by
\[
 \left\{\sum_i\left(\sum_j M_{ij}^2\right)
                    \left(\sum_l N_{il}^2\right)\right\}^{1/2}
 \le\|M\|_F\|N\|_F.
\]
Its local mean has coefficient matrix
$(M^*N)_{jl}$ off the diagonal; excluding $i=j,l$
does not change the sum because $M,N$ have zero diagonal.
The other three shared-index patterns give the same estimate
with a transpose or with the two coordinates exchanged, and
their mean contractions are matrix products of $M,N$ or their
transposes.

Finally, there are two patterns with both indices shared:
$(k,l)=(i,j)$ and $(k,l)=(j,i)$.  Their coefficients are
$M_{ij}N_{ij}$ and $M_{ij}N_{ji}$.  After centering the two
local products, the resulting coefficients are a Hadamard
matrix, its row sums, its column sums, and its total sum.
For instance
\[
 \sum_i\left|\sum_j M_{ij}N_{ij}\right|^2
 \le\sum_i\left(\sum_j M_{ij}^2\right)
            \left(\sum_j N_{ij}^2\right)
 \le\|M\|_F^2\|N\|_F^2,
\]
and the other three estimates follow from the same
Cauchy--Schwarz inequality.  These cases exhaust the possible
index coincidences.  Applying
\eqref{eq:canonical-chaos-bound} to the centered pieces
therefore proves
\begin{equation}
 \|H_AS_B\|_2\le C\|A\|_F\|B\|_F.
 \label{eq:coordinate-function-score-bound}
\end{equation}

It remains to bound $(Bx)\cdot\nabla H_A$.  We give the
contractions separately, since derivatives of a centered
function need not be centered.  For the diagonal part the sum
is
\[
 \sum_{i,k}A_{ii}B_{ik}r'(X_i)X_k.
\]
When $i\ne k$, split $r'$ into its mean and centered part.
The latter has coefficient matrix
$(\diag A)B$, with its diagonal
deleted.  The mean has coefficient vector
$B^*\mathbf a-\mathbf a\mathbin{\odot}\mathbf b$.
Their norms are bounded by, respectively, one and two times
$\|A\|_F\|B\|_F$.  When $i=k$, center the local function
$xr'(x)$, which is in $L^2$ since $r'$ is bounded.  The
coefficient vector and the mean are bounded by the same
product of matrix norms.

Let $M=A_0$.  The derivative of the
off-diagonal part is the sum of
\[
 F(M,B;p',q):=
 \sum_{i\ne j,k}M_{ij}B_{ik}p'(X_i)q(X_j)X_k
\]
and $F(M^*,B;q',p)$.  In $F$, first take $k$ distinct
from $i,j$ and split $p'$ into its mean and centered part.
The centered part has three-index coefficient norm bounded by
\[
 \left\{\sum_i\left(\sum_jM_{ij}^2\right)
                    \left(\sum_kB_{ik}^2\right)\right\}^{1/2}
 \le\|M\|_F\|B\|_F.
\]
The mean part has, off the diagonal, the matrix
$M^*B-M^*\diag B$;
the second term removes $i=k$, and $i=j$ already contributes
zero.  Its Frobenius norm is at most $2\|M\|_F\|B\|_F$.
When $k=i$, center $xp'(x)$; the resulting matrix has entries
$B_{ii}M_{ij}$ and its mean has vector
$M^*\mathbf b$.  When $k=j$, split both
$p'$ and $xq(x)$ into means and centered parts; the coefficients
are the Hadamard matrix $M\mathbin{\odot}B$, its row and column
sums, and its total sum.  The preceding Hadamard estimates
bound all of them.  All the local functions used here belong
to $L^2$, since $p',q$ are bounded and $\E X^2=1$.
The same proof applies to $F(M^*,B;q',p)$.
Using \eqref{eq:canonical-chaos-bound} again proves
\[
 \|(Bx)\cdot\nabla H_A\|_2\le C\|A\|_F\|B\|_F.
\]
Together with \eqref{eq:coordinate-function-score-bound} and
\eqref{eq:coordinate-function-generator-product}, this proves the lemma.
\end{proof}

We can now close the local argument.
The test functions \(H_A\) give a linear functional that sees the
first-order displacement, while the generator bound controls the error made
along the exponential path.
This turns the nonzero rotational Fourier defect into a uniform Hellinger lower bound
for every small matrix perturbation.

\begin{proposition}[Uniform local Hellinger coercivity]
\label{prop:local-fisher-hellinger}
For every non-Gaussian density satisfying
\eqref{eq:fisher-regularity}, there are constants
$c,C,\delta>0$, independent of $n$, such that
\begin{equation}
 c\|B\|_F\le h((I+B)_\#\mu_n,\mu_n)
             \le C\|B\|_F
 \qquad\text{whenever }\|B\|_F\le\delta.
 \label{eq:local-matrix-hellinger}
\end{equation}
The lower bound can be obtained using the fixed finite set of
coordinate frequencies from \cref{lem:fourier-dual-moments}.
\end{proposition}

\begin{proof}
Fix the finite real Fourier sums from \cref{lem:fourier-dual-moments}.
The dependence of \(H_A\) on the direction \(A\) is essential.  Its pairing
with the tangent is quadratic, \(\frac12\|A\|_F^2\), whereas
\(\|H_A\|_2\) is only linear in \(\|A\|_F\); Cauchy--Schwarz will therefore
give a linear Hellinger lower bound.  The generator estimate is used only to
keep the nonlinear remainder cubic, with a coefficient independent of \(n\).

First use exponential coordinates.  Unitarity and the
one-parameter group derivative give
\begin{align*}
 \langle H_A\Omega_n,(U_{\exp A}-I)\Omega_n\rangle
 &=\int_0^1
   \langle U_{\exp(-tA)}(H_A\Omega_n),G_A\Omega_n\rangle\,dt.
\end{align*}
Replacing $U_{\exp(-tA)}(H_A\Omega_n)$ by $H_A\Omega_n$ in the integral
gives the leading term $\tfrac12\E H_AS_A=\tfrac12\|A\|_F^2$.
For a vector in the domain of $G_A$, the fundamental theorem of
calculus gives
$\|(U_{\exp(-tA)}-I)v\|_2\le t\|G_Av\|_2$.
Thus \cref{lem:fourier-generator-bound} and the Fisher upper
bound imply
\[
 \left|\langle H_A\Omega_n,(U_{\exp A}-I)\Omega_n\rangle
                       -\tfrac12\|A\|_F^2\right|
 \le\frac{C_{\mathrm{gen}}\sqrt{C_f}}4\|A\|_F^3.
\]
Cauchy--Schwarz and \eqref{eq:matched-fourier-pairing} now
give, for all sufficiently small $\|A\|_F$,
\[
 h((\exp A)_\#\mu_n,\mu_n)
 \ge\frac{1}{4C_{\mathrm{dual}}}\|A\|_F.
\]
For $\|B\|_{\mathrm{op}}<1/2$, the real power-series logarithm
$A=\log(I+B)$ exists and satisfies
\[
 \|A-B\|_F
 \le\sum_{k\ge2}\frac{\|B\|_{\mathrm{op}}^{k-1}}{k}\|B\|_F
 \le C\|B\|_{\mathrm{op}}\|B\|_F.
\]
Its norm is therefore uniformly comparable to $\|B\|_F$ in a
fixed small chart.  This proves the lower bound for $I+B$;
the upper bound follows from
\cref{lem:finite-hellinger-upper}.
\end{proof}

The half-density pairing supplies the uniform local certificate; tail symmetry
in \cref{sec:necessity} applies it to arbitrary output marginals without a
separate mixed-law variance bound for \(H_A\).

\subsection{Uniform coercivity modulo the exact linear stabilizer}
\label{sec:global-coercivity}

The following global form is proved in \cref{sec:globalization-proof}, after
the infinite-dimensional classification supplies the compactness input; it is
not used in that classification.

\begin{theorem}[Dimension-free global Hellinger coercivity]
\label{thm:global-coercivity}
Assume \eqref{eq:main-assumptions}.  Let \(\mathcal P_n\) denote the allowed
signed-permutation matrices in dimension \(n\), and set
\begin{equation}
 d_n(T)=\min_{Q\in\mathcal P_n}\|T-Q\|_F.
 \label{eq:distance-stabilizer}
\end{equation}
For every \(0<a\le b<\infty\)
there are constants \(c=c(f,a,b)>0\) and \(C=C(f,a,b)<\infty\) such
that, for every \(n\) and every \(T\in\GL(n,\R)\) whose singular
values lie in \([a,b]\),
\begin{equation}
 c\min\{d_n(T)^2,1\}
 \le h(\mu_n,T_\#\mu_n)^2
 \le C\min\{d_n(T)^2,1\}.
 \label{eq:global-coercivity}
\end{equation}
\end{theorem}
